% Daten – bank/daten/ (zusammenbau v0.4, Heft)
\einheitenkopf[t7]{Daten}
\setcounter{aufgabe}{4}

\begin{aufgabe}{Kenngrößen – Prüfungsaufgaben}
\begin{teile}
\teil Weitsprung in m: $4{,}35$; $4{,}62$; $4{,}18$; $4{,}50$; $4{,}27$; $4{,}44$. Spannweite? \hfill (P10 2019 OS) \\ \leerfeld[m]
\teil Messwerte in °C: $6$; $9$; $2$; $11$; $4$; $9$; $7$. Bestimme den Zentralwert (Median). \hfill (P10 2015 OS)
\teil Datenreihe in °C: $16$, $13$, $19$, $14$, $19$, $9$. Bestimme den Median. \hfill (P10 2024 OS)
\teil Werte: $6$; $30$; $45$. Arithmetisches Mittel? \hfill (P10 2016 OS)
\teil Sprunghöhen: $1{,}45$ m; $1{,}50$ m; $1{,}35$ m; $1{,}50$ m. Arithmetisches Mittel? \hfill (P10 2017 OS) \\ \leerfeld[m]
\teil Weiten: $5{,}12$ m; $4{,}96$ m; $5{,}04$ m; $5{,}28$ m. Durchschnittliche Weite? \hfill (P10 2021 OS) \\ \leerfeld[m]
\teil Tageswerte in °C: Mo $17$, Di $19$, Mi $16$, Do $21$, Fr $18$, Sa ?, So $15$. Der Wochendurchschnitt beträgt $18$ °C. Trage den fehlenden Samstagswert ein. \hfill (P10 2022 OS) \\ Sa: \leerfeld °C
\end{teile}
\end{aufgabe}
